\chapter{Design Architecture}\label{chap:design_architecture}

\section{Design Philosophy}

The governing principle of this extension was functionality first, optimisation second. This drove two central decisions. First, floats are represented as 64-bit doubles rather than 32-bit single-precision values, in order to preserve fidelity with Python's native semantics. Second, floats are stored as boxed heap objects rather than as raw values, accepting the cost of indirection in exchange for reusing the compiler's existing tuple and garbage-collection machinery. Performance optimisations were partly explored as part of this work, with the remainder left for future work.

A related decision concerns mixed-type arithmetic. In any operation mixing an integer and a float, the integer operand is implicitly promoted to a float before the operation is performed, matching Python's native behaviour. This keeps a single arithmetic path for floats and confines the type difference to a conversion step at the operands. The design and placement of this conversion are discussed in Section~\ref{sec:coercion}.

\section{Float Representation on the Heap}

The garbage collector uses bit tagging to distinguish pointers, which it must follow, from raw values, which it must leave untouched. Integers are stored shifted one bit to the left, so that their least-significant bit (LSB) is~0, marking the word as a value the collector should not trace. Heap pointers instead carry an LSB of~1, signalling the collector to follow them. Every structure stored on the heap (tuples, closures, and so on) is laid out as a block of words prefixed by a header word.

A double in the IEEE-754 format cannot be tagged in this way, because all 64~bits are needed to encode the value, leaving no bit free for a tag. Worse, the bit that serves as the tag for integers corresponds, in a double, to part of the mantissa; repurposing it would distort the value entirely.

Floats must nonetheless be storable on the heap, since the language permits them as fields of objects, as function arguments, and as elements of tuples. Restricting the language to avoid this was rejected on principle, as the guiding philosophy stated above places functionality over performance. A representation compatible with the existing tagged collector was therefore required.

\subsection{Alternatives}

One option would have been to store floats as raw values directly on the heap. This would require overhauling the garbage collector to move away from bit tagging towards a layout map, that is, a description of the heap that records, for each structure, which words are pointers and which are raw values. Such a scheme would offer a scalable and memory-efficient way to represent any value directly on the heap, and would allow values to be stored inside any other structure without indirection. The drawback is that it conflicts with the architecture of the compiler as built so far, and that constructing and maintaining layout maps would introduce considerable overhead.

The approach I chose instead was to box every float inside a two-word heap block: one header word and one word holding the value of the float. The header identifies the block as a boxed float. This yields a uniform representation that can be used for literals, for variables, and for floats nested in any heap structure, all through the same interface.

The intended behaviour is for the collector to recognise a boxed float from its header and to copy the payload word verbatim, without interpreting those raw bits as a pointer. Section~\ref{sec:gc-extension} describes how this header-aware handling was implemented.

The advantage of boxing is that it offers a uniform design that delivers full functionality with little overhead and fits almost seamlessly into the existing compiler. The disadvantage is the inefficiency that arises from having to reach inside the box on every access to the value, and from allocating a fresh box for every float that is produced, which increases pressure on the collector. These costs were judged acceptable, since the trade-off aligned with the design philosophy of providing functionality before performance.


\section{Conversion of Integers to Floats in Coercion Expressions}\label{sec:coercion}

Python implicitly promotes integers to floats in mixed-type operations---\texttt{3 + 3.14} evaluates to \texttt{6.14}, not a type error. Supporting this behaviour follows directly from the design philosophy of preserving Python's native semantics.

The detection of cross-type operations is performed in the heap allocation pass. This pass already identifies floats to box them, and more generally decides what must happen to each piece of data---whether it needs boxing, unboxing, or conversion. Placing the coercion logic here keeps these semantic decisions in one place.

When the pass detects an operation in which at least one operand is a float and the other is an integer, it wraps the integer operand in a new unary operation, \texttt{int\_to\_float}. Structurally, this operation behaves just like any other unary expression (such as negation); it flows through the subsequent passes unchanged until it reaches instruction selection, where it is lowered to the RISC-V \texttt{fcvt.d.l} instruction that converts a 64-bit integer value to a double-precision float.

\subsection{Alternatives}

One alternative would have been to introduce a dedicated coercion pass before the heap allocation pass. This would separate concerns more cleanly, since one could argue that the purpose of the allocation pass is heap allocation, not type coercion. While technically correct, the trade-off would be introducing an entire new pass for a single transformation in one specific case. This option was rejected as disproportionate; however, should the language require additional type conversions in the future, a unified coercion pass would become architecturally justified.

A second and simpler alternative would have been to handle coercion entirely in instruction selection, using the type dictionary to detect mixed-type operands and emitting the conversion inline. This would confine the change to a single file, making it the least complex solution. The drawback is that instruction selection would take on semantic work---deciding \emph{what} conversion is needed---on top of its existing role of deciding \emph{how} to lower operations to machine instructions. The chosen approach keeps these responsibilities separate: the allocation pass decides that a conversion is needed, and instruction selection decides how to implement it. This separation also makes a future migration to a dedicated coercion pass straightforward, since the conversion is already represented as an explicit operation in the intermediate representation.

\subsection{Possible Improvements}\label{sec:possible-improvements}

Efficiency could be improved by replacing bit tagging with a full layout or type map, which would permit storing raw floats directly on the heap rather than boxing them. Complementarily, the register allocator could be extended to distinguish floats from integers and to maintain a separate region of the stack for floating-point values, so that raw floats could be spilled to the stack just like integers, instead of being boxed and sent to the heap.


\section{Garbage Collector Extension}\label{sec:gc-extension}

Since bit tagging doesn't work for floats - because they need the whole 64 bits - they are boxed. Yet once the Garbage Collector follows the pointer to a boxed float, it still has to decide whether the box's payload is itself a pointer it needs to follow, or just a raw value. To solve this, the header of a boxed float gets a second tag bit (bit 1) marking it as raw data.

Simple as that - if the Garbage Collector reads a header, it checks bit 1, and if it is set, skips over the payload instead of scanning it for pointers, then continues as usual.

The simplicity and seamlessness of this approach is why I decided for it, but there is a more structural reason as well. The collector is a Cheney-style copying collector~\citep{cheney1970nonrecursive}, meaning it scans to-space sequentially rather than by following individual pointers. By the time the scan reaches a given object, it no longer knows which pointer led to it - the object itself is all the scan has to go on. Any alternative that moves the "raw or pointer" information onto the pointer instead of the object breaks that assumption.

\subsection{Alternatives}

One alternative would be to use 2 tag bits on the pointer itself, already encoding what kind of value it points to. This runs into exactly the problem above: the sequential to-space scan never has access to the pointer that led to an object, so the tag would have to be duplicated onto the object anyway, or the sequential scan would have to be replaced with an explicit worklist - giving up the main advantage of the algorithm. It is also harder to extend, since the number of tag bits needed grows with the number of distinct pointer kinds.

Another alternative would be to use two separate address spaces, one for pointers and one for raw values. This would be extensible and clean but is likely overkill for this thesis.

Because of these reasons, the simplicity of encoding the type in the header - which fits how the collector already traverses memory - was selected.


\section{Generalizing the Float Box}

The float box can be generalized to hold an arbitrary number of raw values instead of exactly one, without any change to the front end or to how tuples work.

I kept the box and the tuple as separate concepts rather than merging them. A tuple containing a float is unchanged: each float is still its own length-one box, referenced by a pointer from the tuple's field. What changed is the box-construction mechanism itself, which now takes a length~$n$ - every existing call site still passes $n=1$.

Nothing yet constructs a box with $n>1$; this generalizes the mechanism, not any consumer of it. Reading a boxed float always accesses position~0, which is a fixed offset regardless of $n$, so this required no change.

I did not cap the number of elements a box can hold. The header only spends two of its 64 bits on the forwarding and raw-data flags, leaving roughly 62 bits for the length - far beyond the required $2^{16}$ or $2^{32}$. A cap would add complexity for no benefit, since every value of $n$ is chosen by the compiler itself at compile time.

\subsection{Alternatives}

Representing an "array of floats" as a special-cased tuple - detecting an all-float tuple and storing its elements inline - was rejected, since it conflates two things that are cleanly separate today: a tuple's fields may independently be integers, pointers, or boxed floats, while a box's payload is exclusively one kind of raw data.

A dedicated bit field for the count, separate from the header's existing length field, was also rejected: the length field was never actually limited to one word, only ever used that way. Reusing it keeps the fix backward compatible - for $n=1$ it produces exactly the same header and collector behaviour as before.

\subsection{Possible Improvements}

An array type built on top of this - fixed, compile-time-constant size and indices - would be the natural next step. A genuinely runtime-sized or runtime-indexed array would additionally require computing addresses from a register, since RISC-V's loads and stores only support a base register plus an immediate offset.
